\section{Technical Evaluation}
\label{sec:technical-evaluation}

Behavioral patch review depends on execution evidence that is sufficiently
faithful and stable for a workflow owner to interpret. We therefore evaluated
the evidence pipeline separately from the user study. The evaluation asks
whether the prototype reconstructs observable behavior correctly, whether its
matched comparisons remain stable under the frozen study configuration, and
whether targeted source interventions provide useful discrimination without
being mistaken for proof of causality.

\subsection{Evaluation Questions and Corpus}

We organized the evaluation around three technical questions:

\begin{description}[leftmargin=0pt,labelindent=0pt,style=nextline]
    \item[\textbf{TQ1: Evidence fidelity.}] Do the trace and fact evaluators
    accurately reconstruct the decisions, tool actions, and outcomes relevant
    to each behavioral commitment?
    \item[\textbf{TQ2: Execution stability.}] Do repeated executions of the
    same skill and case support a stable matched comparison under the study
    configuration?
    \item[\textbf{TQ3: Intervention discriminability.}] Do bounded
    interventions expose sensitivity to repair-relevant source regions while
    preserving a visible distinction between source-location evidence and
    whole-candidate validation?
\end{description}

The evaluation corpus contains the complete travel-recovery and
expense-review scenario packs used in the controlled study. Each pack includes
the faulty source skill, an independently authored reference repair, the
product-facing cases, post-decision cases, typed tool schemas, and frozen world
states. Reference behaviors and case-level criteria were specified before the
candidate-generation path was sampled. We exercised the faulty skill,
reference repair, source interventions, and end-to-end candidate path across
all cases using the same model and execution service later used in the study.
Appendix~\ref{app:system-details} reports the exact artifacts, budgets, and
configuration.

\subsection{Measures and Analysis}

For evidence fidelity, two researchers independently inspected a stratified
sample of raw tool traces and terminal world states and annotated the
commitment-relevant facts and event ordering. We compared these annotations
with the system's deterministic evaluator outputs, resolving disagreements
without access to condition labels or candidate provenance. We report exact
agreement, class-specific errors, and inter-rater agreement before resolution.

For execution stability, we measured valid-run rate, modal behavior share, and
agreement of commitment-relevant facts across repetitions of each
skill--case pair. We examined the faulty and reference artifacts separately so
that a stable evaluator could not conceal unstable agent behavior. For matched
comparisons, we additionally verified that original and candidate executions
shared the same task context, tool snapshot, configuration, and side-effect
policy.

For intervention discriminability, we compared the selected source
interventions with interventions on instructions not expected to govern the
target behavior. We measured whether each intervention changed the incident's
commitment-relevant facts, whether it also altered preservation cases, and how
often the selected source region ranked above non-target regions by behavioral
sensitivity. This analysis evaluates the usefulness of the cue under the
scenario corpus; it does not estimate unique causal responsibility among
interacting natural-language instructions.

Finally, we sampled the complete candidate-generation and execution path to
test whether the evidence shown during a realistic repair remained
reconstructable end to end. We report incident success, preservation and
holdout behavior, execution failures, and operational latency. These candidate
samples assess the readiness of the study apparatus rather than the human
benefit of the workspace.

\subsection{Results}

\paragraph{TQ1: Evaluators faithfully reconstructed observable behavior.}
Across \DataNeeded{NUMBER OF MANUALLY AUDITED EXECUTIONS} audited executions,
the deterministic evaluators agreed with the resolved human annotations on
\DataNeeded{OVERALL EXACT AGREEMENT, WITH CONFIDENCE INTERVAL}. Agreement was
\DataNeeded{FACT-CRITERION AGREEMENT} for fact criteria and
\DataNeeded{TRACE-ORDER AGREEMENT} for trace-order criteria. Independent
annotators achieved \DataNeeded{INTERRATER AGREEMENT AND MEASURE} before
resolution. The remaining \DataNeeded{NUMBER} discrepancies involved
\DataNeeded{CHARACTERIZE EXTRACTION OR TRACE-ORDER EDGE CASES}; none
\DataNeeded{CONFIRM WHETHER ANY DISCREPANCY CHANGED A CASE-LEVEL REVIEW STATE}.
We corrected \DataNeeded{REPORT ANY EVALUATOR OR FIXTURE CORRECTIONS MADE
BEFORE FREEZING THE STUDY BUILD} and repeated the affected calibration before
participant data collection.

\paragraph{TQ2: Matched executions were stable enough to expose genuine
uncertainty.}
The frozen corpus produced valid executions in
\DataNeeded{VALID-RUN RATE AND NUMERATOR/DENOMINATOR}. Commitment-relevant
behavior was stable in \DataNeeded{NUMBER OR PROPORTION OF SKILL--CASE PAIRS
MEETING THE PREREGISTERED STABILITY THRESHOLD}, with a median modal share of
\DataNeeded{MEDIAN AND INTERQUARTILE RANGE}. The faulty skills reliably
reproduced the motivating breakdowns, while the reference repairs achieved
their incident targets and preservation criteria in
\DataNeeded{REFERENCE SUCCESS RATES}. The remaining unstable pairs concerned
\DataNeeded{CHARACTERIZE UNSTABLE CASES}; the interface surfaced these as
insufficient evidence rather than collapsing them into a deterministic
outcome. Configuration audits found \DataNeeded{NUMBER} unmatched original--
candidate comparisons after freezing \DataNeeded{EXPECTED TO BE ZERO; REPORT
ANY EXCEPTIONS AND THEIR HANDLING}.

\paragraph{TQ3: Source interventions provided discriminating but bounded
location cues.}
The source regions selected for the two scenario packs changed the
incident-relevant behavior in \DataNeeded{SELECTED-INTERVENTION HIT RATE},
compared with \DataNeeded{NON-TARGET INTERVENTION RATE} for non-target source
regions. They ranked within the top \DataNeeded{RANK OR PERCENTILE} regions by
behavioral sensitivity in \DataNeeded{PROPORTION OF CALIBRATION SAMPLES}.
Their effects were not perfectly local: \DataNeeded{NUMBER OR PROPORTION} also
changed at least one neighboring case, illustrating why the system treats the
intervention as a source-location cue and subsequently executes the complete
candidate across the behavioral scope. No source intervention result was used
as confirmatory evidence for newly generated text.

\paragraph{The complete candidate path supported the planned study tasks.}
Across \DataNeeded{NUMBER OF END-TO-END CANDIDATE SAMPLES}, generated
candidates repaired the motivating incident in
\DataNeeded{INCIDENT SUCCESS RATE}, preserved visible commitments in
\DataNeeded{VISIBLE PRESERVATION RATE}, and preserved generator-withheld
behavior in \DataNeeded{WITHHELD PRESERVATION RATE}. The candidate path also
produced \DataNeeded{NUMBER AND TYPES OF DELIBERATE OR NATURALLY OCCURRING
MISMATCHES} cases in which review should lead to revision, further evidence, or
deferral rather than release. Median end-to-end latency was
\DataNeeded{LATENCY WITH DISTRIBUTION}, and execution or service failures
occurred in \DataNeeded{FAILURE RATE}. These results established that both
study conditions could expose the same substantive repair evidence within the
planned task duration.

\subsection{Summary and Evidentiary Boundaries}

The technical evaluation supports the use of deterministic trace and fact
evidence for the two controlled domains and shows that the prototype can
distinguish stable behavior from insufficient evidence. It also supports the
narrower role assigned to source interventions: they can direct attention to
behaviorally relevant text, but their effects may extend across cases and do
not identify a unique causal instruction. The evaluation does not establish
reliability for arbitrary skills, tools, models, or deployment environments.
The controlled user study therefore tests whether people benefit from this
evidence within the evaluated scope, not whether \systemname{} certifies a
revision for unrestricted use.
